\documentclass[10pt]{article}
% =====================================================================
%  Préambule commun aux CV FR et EN (style RenderCV / sb2nov)
%  Modifier ici la mise en page : elle s'applique aux deux langues.
% =====================================================================
\usepackage[
    a4paper,
    ignoreheadfoot,
    top=0.8 cm,
    bottom=0.8 cm,
    left=1.8 cm,
    right=1.8 cm,
    footskip=0.5 cm,
]{geometry}
\usepackage{titlesec}
\usepackage{tabularx}
\usepackage{array}
\usepackage[dvipsnames]{xcolor}
\definecolor{primaryColor}{RGB}{0, 79, 144}
\usepackage{enumitem}
\usepackage{fontawesome5}
\usepackage{amsmath}
\usepackage[
    pdftitle={CV Philibert Pappens},
    pdfauthor={Philibert Pappens},
    colorlinks=true,
    urlcolor=primaryColor
]{hyperref}
\usepackage{calc}
\usepackage{bookmark}
\usepackage{changepage}
\usepackage{paracol}
\usepackage{ifthen}
\usepackage{needspace}
\usepackage{iftex}

% PDF lisible par les ATS
\ifPDFTeX
    \input{glyphtounicode}
    \pdfgentounicode=1
    \usepackage[utf8]{inputenc}
    \usepackage{lmodern}
\fi

\AtBeginEnvironment{adjustwidth}{\partopsep0pt}
\pagestyle{empty}
\setcounter{secnumdepth}{0}
\setlength{\parindent}{0pt}
\setlength{\topskip}{0pt}
\setlength{\columnsep}{0cm}

\titleformat{\section}{\needspace{4\baselineskip}\bfseries\large}{}{0pt}{}[\vspace{0.5pt}\titlerule]
\titlespacing{\section}{-1pt}{0.15 cm}{0.07 cm}

\renewcommand\labelitemi{$\circ$}
\newenvironment{highlights}{
    \begin{itemize}[topsep=0.07 cm, parsep=0.07 cm, partopsep=0pt, itemsep=0pt, leftmargin=0.3 cm + 10pt]
}{
    \end{itemize}
}

\newenvironment{onecolentry}{
    \begin{adjustwidth}{0.2 cm + 0.00001 cm}{0.2 cm + 0.00001 cm}
}{
    \end{adjustwidth}
}

\newenvironment{twocolentry}[2][]{
    \onecolentry
    \def\secondColumn{#2}
    \setcolumnwidth{\fill, 4.5 cm}
    \begin{paracol}{2}
}{
    \switchcolumn \raggedleft \secondColumn
    \end{paracol}
    \endonecolentry
}

\newenvironment{header}{
    \setlength{\topsep}{0pt}\par\kern\topsep\centering\linespread{1.5}
}{
    \par\kern\topsep
}

% Raccourcis
\newcommand{\entrysep}{\vspace{0.1 cm}}   % espace entre deux entrées
\newcommand{\bulletsep}{\vspace{0.07 cm}} % espace titre -> puces
\newcommand{\cpp}{C\texorpdfstring{$++$}{++}}

% Liens externes avec flèche
\let\hrefWithoutArrow\href
\renewcommand{\href}[2]{\hrefWithoutArrow{#1}{\ifthenelse{\equal{#2}{}}{ }{#2 }\raisebox{.15ex}{\footnotesize \faExternalLink*}}}

% En-tête (identique FR/EN, seul le lieu change)
\newcommand{\cvheader}[2]{% #1 = lieu, #2 = libellé du lien vers le site
    \begin{header}
        \textbf{\fontsize{24 pt}{24 pt}\selectfont Philibert Pappens}

        \vspace{0.2 cm}

        \small
        \mbox{{\color{black}\footnotesize\faMapMarker*}\hspace*{0.1cm}#1}%
        \hspace{0.25 cm plus 0.1 cm}%
        \mbox{\hrefWithoutArrow{mailto:philibert.pappens@gmail.com}{\color{black}{\footnotesize\faEnvelope[regular]}\hspace*{0.1cm}philibert.pappens@gmail.com}}%
        \hspace{0.25 cm plus 0.1 cm}%
        \mbox{\hrefWithoutArrow{tel:+33783458038}{\color{black}{\footnotesize\faPhone*}\hspace*{0.1cm}+33 (0)7 83 45 80 38}}%
        \hspace{0.25 cm plus 0.1 cm}%
        \mbox{\hrefWithoutArrow{https://linkedin.com/in/philibert-pappens-993468313}{\color{black}{\footnotesize\faLinkedinIn}\hspace*{0.1cm}Philibert Pappens}}%
        \hspace{0.25 cm plus 0.1 cm}%
        \mbox{\hrefWithoutArrow{https://philibertpap.github.io/site_perso}{\color{black}{\footnotesize\faGlobe}\hspace*{0.1cm}#2}}%
    \end{header}
}


\begin{document}

\cvheader{Stockholm, Suède}{Site web}

\begin{center}
    \small
    Élève ingénieur de l'École polytechnique en double diplôme à KTH (architecture navale), spécialisé en mécanique des matériaux et des structures. Expériences en recherche appliquée sur les composites, en industrie et à l'international ; j'aime découvrir vite de nouveaux sujets et relier l'analyse technique aux enjeux d'un projet.
\end{center}


\section{Formation}

    \begin{twocolentry}{
    \textit{Stockholm, Suède}

    \textit{Août 2026 – Déc. 2027}}
        \textbf{KTH Royal Institute of Technology}

        \textit{Master of Science in Engineering -- Naval Architecture (double diplôme)}
    \end{twocolentry}

    \bulletsep
    \begin{onecolentry}
        \begin{highlights}
            \item \textbf{Cours :} structures marines, conception de navires, fatigue des structures soudées, structures légères et éléments finis, composites, hydromécanique
            \item Projet de conception d'un navire roulier (lac Vänern -- Ruhr) -- mémoire de master à l'automne 2027
        \end{highlights}
    \end{onecolentry}

    \entrysep

    \begin{twocolentry}{
    \textit{Palaiseau, Essonne}

    \textit{Sept. 2023 – Aujourd'hui}}
        \textbf{École polytechnique, IP Paris}

        \textit{Diplôme d'ingénieur polytechnicien}
    \end{twocolentry}

    \bulletsep
    \begin{onecolentry}
        \begin{highlights}
            \item \textbf{Spécialisation mécanique des matériaux et des structures} (3A) -- 4A en double diplôme à KTH
        \end{highlights}
    \end{onecolentry}

    \entrysep

    \begin{twocolentry}{
    \textit{Versailles, Yvelines}

    \textit{Sept. 2021 – Juil. 2023}}
        \textbf{Lycée privé Sainte-Geneviève}

        \textit{Classes préparatoires aux grandes écoles -- MP, option informatique}
    \end{twocolentry}


\section{Expérience professionnelle}

    \begin{twocolentry}{
    \textit{Lorient, Morbihan}

    \textit{Mars 2026 – Juil. 2026}}
        \textbf{Institut de recherche Dupuy de Lôme (IRDL) -- Stage de recherche}

        \textit{Équipe DEMAT (comportement et durabilité des matériaux)}
    \end{twocolentry}

    \bulletsep
    \begin{onecolentry}
        \begin{highlights}
            \item Implémentation des conditions de Bloch--Floquet dans Abaqus (scripting Python) pour prédire la résistance en compression (microflambage) de stratifiés carbone à plis unidirectionnels, représentatifs des hydrofoils de voile de compétition
            \item Conduite autonome du sujet (bibliographie, choix de méthode, validation par comparaison à des résultats numériques antérieurs obtenus par d'autres méthodes) -- outil intégré au workflow de l'équipe
        \end{highlights}
    \end{onecolentry}

    \entrysep

    \begin{twocolentry}{
    \textit{Bonn, Allemagne}

    \textit{Juin 2025 – Sept. 2025}}
        \textbf{Deutsche Telekom Technik -- Stage de découverte de l'entreprise}

        \textit{Équipe Digital Transformation -- Fiber Factory}
    \end{twocolentry}

    \bulletsep
    \begin{onecolentry}
        \begin{highlights}
            \item Participation à des projets de transformation numérique et d'amélioration des processus ; rédaction d'un playbook interne pour l'équipe
            \item Travail quotidien en environnement entièrement germanophone au sein d'un groupe international
        \end{highlights}
    \end{onecolentry}

    \entrysep

    \begin{twocolentry}{
    \textit{Guingamp}

    \textit{Oct. 2023 – Avril 2024}}
        \textbf{Officier de gendarmerie -- Stage de formation humaine et militaire}

        \textit{Gendarmerie nationale -- Formation à l'AMGN (Melun)}
    \end{twocolentry}

    \bulletsep
    \begin{onecolentry}
        \begin{highlights}
            \item Participation à des missions de sécurité publique : patrouilles, enquêtes, maintien de l'ordre
            \item Suivi du commandement lors d'événements sensibles ; rigueur et communication sous pression
        \end{highlights}
    \end{onecolentry}


\section{Autres expériences et projets}

    \begin{twocolentry}{
    \textit{Paris}

    \textit{2023 – Août 2026}}
        \textbf{Assistant chef de troupe}

        \textit{Scouts unitaires de France -- Scouts marins}
    \end{twocolentry}

    \bulletsep
    \begin{onecolentry}
        \begin{highlights}
            \item Encadrement d'un groupe de 20 adolescents -- deux camps d'été de deux semaines avec 4--5 jours de navigation -- CEP1 -- infirmier de la troupe (PSC1)
        \end{highlights}
    \end{onecolentry}

    \entrysep

    \begin{twocolentry}{

    \textit{Sept. 2024 – Mars 2026}}
        \textbf{Projets de modélisation et de programmation}

        \textit{École polytechnique}
    \end{twocolentry}

    \bulletsep
    \begin{onecolentry}
        \begin{highlights}
            \item Gréement de voilier : calculs analytiques et éléments finis sous Cast3M (flambement, fatigue, modes propres)
            \item Simulation de l'impact iceberg/Titanic : modèle éléments finis de la coque rivetée sous Python/FEniCS
            \item Informatique quantique : problème d'isomorphisme de graphes par calcul adiabatique (Qiskit)
            \item Application web (HTML, CSS, JavaScript, PHP, SQL) et version 3D du jeu TRON en \cpp
        \end{highlights}
    \end{onecolentry}


\section{Compétences}

    \begin{onecolentry}
        \textbf{Langues :} français (langue maternelle), allemand (C1), anglais (C1+, Linguaskill), suédois (débutant)
    \end{onecolentry}

    \bulletsep

    \begin{onecolentry}
        \textbf{Programmation :} Python, MATLAB, \cpp, Java, OCaml, HTML/CSS/JavaScript, SQL/PHP, \LaTeX \quad $|$ \quad \textbf{Simulation :} Abaqus, Cast3M, FEniCS, Qiskit
    \end{onecolentry}

    \bulletsep

    \begin{onecolentry}
        \textbf{Passe-temps :} gardien de but (équipe de football de l'école) ; responsable communication et chanteur -- \textit{Ensemble vocal de l'École polytechnique} ; orgue (8 ans, dont 3 en conservatoire) et piano
    \end{onecolentry}

\end{document}
